\section{第~\ref{sec:dofa}~章：借助\nt{DOPA}学习}

突触的机制（量子、符合检测器、钙、可塑性窗口）以及多巴胺作为预测误差的行为是直接测量的；规则导向梯度以及广播的代价是定理；选择学习的结果是按本书模型的计算；计算误差的电路是假说。

{\footnotesize
\setlength{\tabcolsep}{3pt}\renewcommand{\arraystretch}{1.15}
\begin{longtable}{>{\raggedright}p{0.5cm}>{\raggedright}p{4.0cm}>{\raggedright}p{5.6cm}>{\raggedright}p{2.3cm}>{\raggedright\arraybackslash}p{1.7cm}}
\toprule
序 & 命题 & 实验与测量内容 & 来源 & 状态 \\
\midrule\endfirsthead
\toprule
序 & 命题 & 实验与测量内容 & 来源 & 状态 \\
\midrule\endhead
1 & 在脑中没有发现人工网络那种形式的误差反向传播 & 没有直接实验：这是一个关于“不存在”的命题。综述分析了该算法需要什么（与前向连接对应的反向连接、单独的阶段），以及已提出的哪些生物学近似。 & \cite{Lillicrap2020, WhittingtonBogacz2019} & 假说 \\
2 & 权重分解为释放位点数、释放概率和对一份递质的响应：$w=N_s\,p\,A_1$~\eqref{eq:quantal} & 释放减少时的蛙神经肌肉突触：接收者的响应呈台阶式波动，每级是自发的单份微小响应的整数倍；每个脉冲释放的份数服从泊松分布。 & \cite{delCastillo1954} & 已测量 \\
3 & 接收者的受体是符合检测器；钙启动权重变化 & 小鼠神经元的膜片钳：在静息电位下，Mg$^{2+}$离子堵塞由谷氨酸打开的通道，去极化时离开。海马脑片：接收者中的钙螯合剂阻断长时程增强，而用光在细胞内释放钙本身就能增强传递。 & \cite{Nowak1984, Malenka1988calcium} & 已测量 \\
4 & 决定可由任何一侧执行：接收者一侧$A_1$增大，发送者一侧$p$改变 & 在单个树突棘上方用光释放谷氨酸，引起树突棘及其受体电流迅速而持久地增长；增强伴随着 AMPA 受体的嵌入。接收者去极化会暂时抑制发送者的释放；这一效应由逆向穿过间隙的内源性大麻素传递（在\nt{GABA}输入上证明）。短时耗竭取决于发送者的释放概率。 & \cite{Matsuzaki2004, Malinow2002, WilsonNicoll2001, Tsodyks1997} & 已测量 \\
5 & 发送者和接收者同时活跃时突触增强~\eqref{eq:hebb} & 麻醉兔：沿海马输入通路施加短串脉冲后，响应增强可持续$30$~min至$10$~h。在海马脑片中，接收者的脉冲只增强同一时刻活跃的那些突触。 & \cite{Bliss1973LTP, Kelso1986hebb} & 已测量 \\
6 & 规则~\eqref{eq:oja}找到输入相关的主特征向量（命题~\ref{prop:oja}） & 定理；证明见本章。没有神经元学会其输入主成分的实验；突触中限制权重增长的机制尚未确定。 & \cite{Oja1982} & 理论 \\
7 & 按时间的赫布规则：窗口约$\pm20\ms$（图~\ref{fig:stdp}）；重复的序列中会长出链 & 培养的海马神经元对：接收者的脉冲在发送者脉冲之后$20\ms$内产生持久增强，之前$20\ms$内产生减弱；二者都依赖 NMDA 受体。图中$A_-=0.6A_+$的比例只是示意。网络中如此长出链，没有直接测量。 & \cite{BiPoo1998} & 已测量 \\
8 & 资格迹~\eqref{eq:threefactor}：共同活动后$1$--$2\,\text{s}$到达的多巴胺增强连接，更晚则不然（命题~\ref{prop:threefactor}） & 纹状体脑片，分别光遗传刺激谷氨酸和多巴胺输入：只有当多巴胺在谷氨酸输入后$0.3$--$2\,\text{s}$到达时，树突棘才长大；窗口由 cAMP 的快速上升和分解决定。在皮层中，配对刺激留下分别用于增强和减弱的迹，单胺受体在秒级窗口内把它们转化为长期变化。 & \cite{Yagishita2014, He2015eligibility} & 已测量 \\
9 & 没有多巴胺也有长达数秒的窗口：第三个信号是接收者的强烈兴奋 & 活体小鼠，CA1 区：接收者中的一次平台事件增强在其前后数秒内到达的输入，一次经过就形成位置野。 & \cite{Bittner2017} & 已测量 \\
10 & 三因子规则的平均步长沿奖励梯度~\eqref{eq:perturbation}（命题~\ref{prop:gradient}） & 基于随机偏离的学习的梯度定理；本章在小噪声下推导。 & \cite{Williams1992} & 理论 \\
11 & 噪声即搜索：网络从自身的随机偏离中学习 & 幼年斑胸草雀：关闭基底神经节回路的输出核，会急剧且可逆地消除鸣唱的变异性。成年斑胸草雀：在音高偏向平均值一侧的音节演唱上施加干扰，鸟会迅速把音高移向另一侧——从自身的离散中学习。 & \cite{Olveczky2005, TumerBrainard2007} & 已测量 \\
12 & 当$\tau_e=1\,\text{s}$、$\delta=R-\bar R$时能学会二选一；$\tau_e=50\ms$时学不会；不减去期望时权重无限增长，学习速率大时网络可能固定在更差的选择上 & \texttt{models/dofa\_learning.py}，$\eta=0.05$，$500$次尝试，$6$个随机种子。$\tau_e=1\,\text{s}$：选择$A$的比例$0.65$--$0.79\to0.96$--$1.00$，权重$0.85$--$1.33$。$\tau_e=50\ms$：$0.38$--$0.55$，权重不变。$\delta=R$：获胜者权重$860$--$1190$。$\delta=R$，$\eta=0.5$，$3$个种子：一个固定在$B$上（$0.04\to0$）。脚本输出全部四种情况；重算结果与本章一致。 & \texttt{models/\allowbreak dofa\_\allowbreak learning.py} & 模型 \\
13 & 依靠一个公共信号学习，所需时间与噪声神经元数目成正比地增长（命题~\ref{prop:broadcastcost}） & 线性网络的学习曲线：扰动神经元比精确梯度慢的倍数等于输出神经元数，扰动权重还要再慢输入数那么多倍。 & \cite{Werfel2005} & 理论 \\
14 & \nt{DOPA}的输出主要通向动作选择区域 & 大鼠单个黑质纹状体\nt{DOPA}神经元轴突的完整重建：一个细胞的分支覆盖纹状体体积的$0.45$--$5.7\,\%$（平均$2.7\,\%$），纹状体外几乎没有分支。 & \cite{Matsuda2009} & 已测量 \\
15 & 皮层学习预测自己的输入，误差在原地计算 & 综述：感觉皮层中存在对预期输入与实际输入失配的响应。这是否为皮层的普遍学习规则，尚未证明。 & \cite{Keller2018} & 假说 \\
16 & 多巴胺的行为像误差~\eqref{eq:ctd}：尖峰、向预示信号的转移、低谷（图~\ref{fig:ctdcases}） & 猴子的\nt{DOPA}神经元：对意外的果汁出现尖峰；学习后对预示信号出现尖峰，对承诺的果汁无响应；果汁未到时出现低谷。连续形式~\eqref{eq:ctd}是理论。 & \cite{Schultz1997, Doya2000} & 已测量 \\
17 & 计算$\dd V/\dd t$并减去期望的电路 & 小鼠腹侧被盖区的记录和光遗传学：\nt{DOPA}神经元从奖励中减去期望；邻近的\nt{GABA}神经元在奖励可以预期时抑制它们，兴奋或抑制这些神经元会改变误差。导数如何计算，尚未证明。 & \cite{Eshel2015arithmetic} & 假说 \\
18 & \nt{DOPA}神经元是起搏器；低谷比尖峰浅 & 在脑片中\nt{DOPA}神经元在缓慢的起搏电位上自发地、非常规则地放电。在猴子中，频率定量地跟随正误差，而负误差只得到微弱的传递。 & \cite{GraceOnn1989, Bayer2005} & 已测量 \\
19 & 行动者与评论器（图~\ref{fig:actorcritic}） & 算法来自强化学习理论。在人类（fMRI）中，腹侧纹状体无论是否需要选择都表现得像评论器，背侧纹状体只在需要选择动作时如此。 & \cite{SuttonBarto2018, ODoherty2004actor} & 假说 \\
20 & 在带第二族受体的神经元中，规则里$\delta$的符号是反的 & 纹状体脑片：在带 D1 受体的神经元中，配对刺激产生需要 D1 的增强；在带 D2 的神经元中产生需要 D2 的减弱。 & \cite{Shen2008} & 已测量 \\
\bottomrule
\end{longtable}}
